% =====================================================================
%  3 Data and Experimental Setup
%  Sources: CLAUDE.md section 11 (design details), S/master_grid.py
%  (prompts, conditions, spans), S/select_items.py (selection),
%  reports/v3_plan.md (decisions, success criterion),
%  reports/v3_step1_explained.md (attention and steering explained),
%  reports/v2_behavioural_analysis.md section 1 (statistical design).
%  S/ = D3code-main/scripts/
% =====================================================================
\section{Data and Experimental Setup}
\label{sec:method}

% Short roadmap: data -> models and scoring -> items -> prompts and
% conditions -> attention measures -> profiling -> steering -> evaluation.
\fillin{One-paragraph overview of the pipeline (maybe a pipeline figure).}

% ---------------------------------------------------------------------
\subsection{Dataset}
\label{sec:method:data}
% - D3CODE (Davani et al. 2024): texts, raters, regions, countries,
%   5-point Likert scale -> binarised how? (check: we use "offensive" =
%   rating >= ? ; see S/analysis_v2_behaviour.py, D3_RATINGS_CSV and the
%   group share computation).
% - Rater attributes: region (8), gender, age, SES, MFQ-2 moral
%   foundations (6).
% - Numbers below: CLAUDE.md section 9 "Data".
\fillin{Describe D3CODE, the rating scale, how a rating is binarised, the rater attributes.}
\src{CLAUDE.md §9 Data}

\begin{table}[t]
  \centering
  \caption{\dthree{} at a glance, and the 500 texts used in the main experiments.}
  \label{tab:data}
  \begin{tabular}{lrr}
    \toprule
                                        & all \dthree{}   & our 500 texts \\
    \midrule
    texts (unique, $\geq 10$ raters)    & \num{4519}      & 500 \\
    raters                              & \num{4309}      & \fillin{n} \\
    ratings                             & $\approx$\num{145000} & \fillin{n} \\
    cultural regions                    & 8               & 8 \\
    share of ratings ``offensive''      & 0.32            & \fillin{x} \\
    mean offensive share per text       & 0.34            & 0.21 \\
    texts majority-offensive            & 21.7\%          & 13 texts \\
    mean rating entropy                 & 1.77            & 1.49 \\
    \bottomrule
  \end{tabular}
\end{table}

% ---------------------------------------------------------------------
\subsection{Models and scoring}
\label{sec:method:models}
% - \llama{} and \qwen{}, chat templates, bf16 weights, single forward
%   pass, no generation.
% - P(yes) from the two answer-token logits, computed in fp32 from the
%   final hidden state (DecisionHead): bf16 output logits quantise P(yes)
%   in steps of 0.06 to 0.125 log-odds (CLAUDE.md section 13).
% - Shift = log-odds(framed) - log-odds(unframed), per text.
We evaluate \llama{} \citep{grattafiori2024llama} and \qwen{} \citep{qwen2025qwen25}.
\fillin{Prompting setup, chat template, single forward pass.}
For a prompt $x$, let $z_0$ and $z_1$ be the logits of the answer tokens
``0'' and ``1'' at the last position. We define
\begin{equation}
  \pyes = \frac{\exp z_1}{\exp z_0 + \exp z_1},
  \qquad
  \ell(x) = \log\frac{\pyes}{1-\pyes} = z_1 - z_0 ,
  \label{eq:pyes}
\end{equation}
and measure the effect of a framing $c$ on text $t$ as the shift
$\Delta_c(t) = \ell(x_c(t)) - \ell(x_\varnothing(t))$ relative to the unframed prompt.
\fillin{Why log-odds and not probabilities (no ceiling compression).}

% ---------------------------------------------------------------------
\subsection{Item selection}
\label{sec:method:items}
% - All 4,519 texts scored with both models (S/select_items.py score).
% - Band: unframed P(yes) in [0.02, 0.98] for BOTH models (ceiling effect
%   found in the pilots, appendix A). 554 texts qualify.
% - 500 drawn with category quotas: 153 social-group (all available),
%   275 random, 72 moral. Split A (selection, 250) / B (held out, 250),
%   stratified by category. Frozen list: S/items/items_v2.csv.
% - Same items for both models (equal-conditions rule).
\fillin{Selection procedure and its rationale.}
\src{S/select\_items.py}

% ---------------------------------------------------------------------
\subsection{Prompt formats}
\label{sec:method:prompts}
% Full texts of all 10 formats in appendix B. Explain the three families:
% perspective (prefix, reworded, imagine, consider, context, context_long),
% position (infix, suffix), and question type (embedded vs target:
% "as X, is it offensive?" vs "is it offensive for X?").
The unframed prompt is
\begin{promptbox}
Text: "\{text\}"\\
Is this text offensive? Answer with only 0 (no) or 1 (yes).
\end{promptbox}
\Cref{tab:variants} lists how a framing is inserted.
\fillin{Rationale for each family of formats.}

\begin{table}[t]
  \centering
  \small
  \caption{Prompt formats. \texttt{X} is the condition phrase (\cref{tab:conditions}),
    \texttt{<base>} the unframed prompt, \texttt{<instr>} its question line.}
  \label{tab:variants}
  \begin{tabularx}{\linewidth}{l>{\ttfamily\raggedright\arraybackslash}X}
    \toprule
    format & \normalfont template \\
    \midrule
    \variant{prefix}       & You are X. <base> \\
    \variant{reworded}     & Assume you are X. <base> \\
    \variant{imagine}      & Imagine you are X. <base> \\
    \variant{consider}     & Consider that you are X. <base> \\
    \variant{context}      & You are X. Your background and everyday experience shape how you read things. <base> \\
    \variant{context\_long}& \ldots{} as context + Answer from that perspective. <base> \\
    \variant{infix}        & Text: "\ldots" / You are X. <instr> \\
    \variant{suffix}       & <base> You are X. \\
    \variant{embedded}     & Text: "\ldots" / As X, is this text offensive? \ldots \\
    \variant{target}       & Text: "\ldots" / Do you think this text is offensive for X? \ldots \\
    \bottomrule
  \end{tabularx}
\end{table}

% ---------------------------------------------------------------------
\subsection{Conditions and controls}
\label{sec:method:conditions}
% - 14 demographic phrases, 8 content-free placebos (hobbies and habits,
%   nothing a rater attribute could encode), 2 per noun length.
% - Exact token-length matching across both tokenizers in all formats
%   (S/check_matching.py).
% - Place-name controls (run B): same place tokens, no belonging.
% - Full factorial: 22 conditions x 10 formats x 500 texts = 110,000
%   framed prompts + 500 unframed per model.
\fillin{Conditions, placebo design principles, length matching, place controls.}

\begin{table}[t]
  \centering
  \small
  \caption{Condition phrases. Placebos come in pairs per noun length; place-name
    controls keep the region's tokens but drop the belonging.}
  \label{tab:conditions}
  \begin{tabularx}{\linewidth}{lX}
    \toprule
    axis & phrases \\
    \midrule
    region (8)   & someone from Western Europe / North America / Latin America / the Arab Culture region /
                   the Indian Cultural Sphere / the Sinosphere / Sub Saharan Africa / Oceania \\
    gender (2)   & a man, a woman \\
    religion (2) & a Christian, a Muslim \\
    identity (2) & an Arab, a European \\
    \midrule
    placebo (8)  & a runner, a cyclist, an early riser, someone who plays chess, someone who collects
                   old maps, someone who enjoys crossword puzzles, someone who usually takes the stairs,
                   someone who prefers tea over coffee \\
    \midrule
    place (8)    & someone researching / mapping $\langle$region$\rangle$, for Western Europe, Latin America,
                   the Arab Culture region, Sub Saharan Africa \\
    \bottomrule
  \end{tabularx}
\end{table}

% ---------------------------------------------------------------------
\subsection{Attention measures}
\label{sec:method:attention}
% - Attention of the last (decision) position to four segments: identity
%   phrase, instruction, item text, other (incl. sink and template).
% - Sink heads flagged and excluded (Xiao et al. 2024).
% - Value-weighted attention (Kobayashi et al. 2020): alpha_ij * ||v_j W_O||,
%   computed via the Gram matrix of W_O^h (S/master_grid.py ValueNorms).
% - Per token (divide by span length) because phrase length drives raw
%   attention.
For head $h$ in layer $l$, let $\alpha^{l,h}_{ij}$ be the attention weight from
position $i$ to $j$. The share of segment $S$ at the decision position $n$ is
\begin{equation}
  A^{l,h}(S) = \sum_{j \in S} \alpha^{l,h}_{nj},
  \qquad
  \tilde A^{l,h}(S) = \frac{\sum_{j\in S} \alpha^{l,h}_{nj}\,\lVert v^{l,h}_j W^{l,h}_O\rVert}
                           {\sum_{j} \alpha^{l,h}_{nj}\,\lVert v^{l,h}_j W^{l,h}_O\rVert},
  \label{eq:attention}
\end{equation}
where $\tilde A$ is the value-weighted variant \citep{kobayashi2020attention}.
\fillin{Sink handling, per-token normalisation, the correlational candidate rule (step 1).}

% ---------------------------------------------------------------------
\subsection{Causal head profiling}
\label{sec:method:profiling}
% - PASTA single-head steering (Zhang et al. 2024), one head at a time,
%   on 48 texts of split A, prefix / infix / target.
% - Head effect = change in log-odds under steering, region vs placebo
%   and region vs place-name controls ("place_specific").
% - Two nulls: placebo-pair (stairs vs tea) and place-pair (researching
%   vs mapping); head passes if beyond the 99th percentile of both.
% - Dampeners (steering lowers region answer) vs amplifiers; frozen
%   before split B (S/head_selection_v3.py, results/v3/analysis/heads_rule.txt).
\fillin{Profiling procedure, nulls, selection rule.}
\src{reports/v3\_runB\_results.md §2}

% ---------------------------------------------------------------------
\subsection{Steering}
\label{sec:method:steering}
% - Attention steering (PASTA): in the selected heads, at every query
%   position, multiply attention to the span by strength k and
%   renormalise. Strengths 0, 0.5, 2, 4, 10, 100 (k = 100 means 0.01 on
%   everything else in PASTA's parameterisation).
% - Arms: dampeners, amplifiers (identity span), text steering (item
%   text, top 10 heads from text profiling), activation steering
%   (mean framed - unframed hidden state, middle layer 16 / 14, answer
%   position), random head sets with the same layer mix (3 per setting).
% - Upper bounds: remove identity attention in all heads at the answer
%   position only vs at all positions.
% - All on split B, prefix / infix / target, 4 regions + gender +
%   religion + placebos + place controls + unframed.
In the steered heads $H$, attention weights to the span $S$ are rescaled at every
query position $i$ and renormalised:
\begin{equation}
  \hat\alpha^{l,h}_{ij} = \frac{k^{\mathbb 1[j\in S]}\,\alpha^{l,h}_{ij}}
                               {\sum_{j'} k^{\mathbb 1[j'\in S]}\,\alpha^{l,h}_{ij'}},
  \qquad (l,h)\in H,
  \label{eq:pasta}
\end{equation}
with strength $k$ \citep{zhang2024pasta}. For activation steering we add
$\lambda\,\bar{\bm d}$ to the residual stream at layer $l^\ast$ and the answer position, where
$\bar{\bm d}$ is the mean difference between framed and unframed hidden states
\citep{turner2023activation}.
\fillin{Check that \cref{eq:pasta} matches S/ablation\_heads.py exactly (multiplying the span vs dividing the rest is equivalent after renormalisation); describe arms and controls.}

% ---------------------------------------------------------------------
\subsection{Evaluation and statistics}
\label{sec:method:eval}
% - Phrase as the statistical unit for content claims (phrases are the
%   same on every item, item-level p-values are too optimistic).
% - Length-adjusted placebo null: regression of placebo shifts on phrase
%   length, demographic compared to the line.
% - Partial Spearman controlling for item-text length and baseline.
% - BH-FDR where many tests.
% - Distance to real raters (success criterion, supervisor 10 Oct):
%   D = mean_t |P(yes) - s_g(t)| on split B, texts with >= 4 group raters.
% - Constant-shift baseline: one log-odds offset fitted on split A.
% - Standard metrics: precision, recall, F1, macro-F1, ROC-AUC, Brier.
The success criterion for steering is the mean absolute distance to the real group
\begin{equation}
  D_g = \frac{1}{|T_g|}\sum_{t\in T_g}\bigl|\,\pyes(x_g(t)) - s_g(t)\bigr| ,
  \label{eq:distance}
\end{equation}
where $s_g(t)$ is the share of group $g$'s raters who rate $t$ as offensive and
$T_g$ are the held-out texts with at least four raters from $g$.
\fillin{Statistical units, nulls, the constant-shift baseline, standard metrics.}

% ---------------------------------------------------------------------
\subsection{Pilot studies}
\label{sec:method:pilots}
% Short: what v1 found and which design choice each lesson led to
% (details in appendix A):
%   ceiling effect -> baseline band; one odd placebo -> placebo pool;
%   phrase length -> exact token matching; unequal item sets -> shared
%   list; bf16 rounding -> fp32 decision logits; span bug -> search from
%   the end.
\fillin{Two or three sentences, pointer to \cref{app:pilots}.}
